\subsection{变系数渐近展开}

可以按如下方式推广主部的概念以及渐近展开的概念。设 E 是一个由实（相应地，复）函数组成的比较尺度，对其中每个函数，存在 F 的一个集使函数在其上任何点都不为零。此外，设 C 是 \(\mathcal{H}(\mathfrak{F}, V)\) 中的函数所成的一个集，满足下列三个条件：

\((\mathrm{CO}_{\mathrm{I}})\) 对每个函数 \(\mathbf{a} \in \mathcal{C}\) 有 \(\mathbf{a} \preccurlyeq 1\) 。

\((\mathrm{CO}_{\mathrm{II}})\) 关系 \(\mathbf{a} \ll 1\) 对函数 \(\mathbf{a} \in \mathcal{C}\) 而言蕴含 \(\mathbf{a} = 0\) 。

\((\mathrm{CO}_{\mathrm{III}})\mathcal{C}\) 是 R（相应地， C）上的一个向量空间。

现在设 \(\mathbf{f}\) 是 \(\mathcal{H}(\mathfrak{F}, V)\) 中的任一函数；如果存在一个函数 \(g \in \mathcal{E}\) 和一个非零函数 \(\mathbf{a} \in \mathcal{C}\) 使得 \(\mathbf{f} - \mathbf{a}g \ll g\) ，就说 \(\mathbf{a}g\) 是 \(\mathbf{f}\) 的一个主部（相对于比较尺度 \(\mathcal{E}\) 与系数域 \(\mathcal{C}\) ）。如果这样的主部存在，则它是唯一的：因为设有两个主部 \(\mathbf{a}_1g_1\) 与 \(\mathbf{a}_2g_2\) ；不可能有 \(g_1 \ll g_2\) ，因为由 \((\mathrm{CO_I})\) 将可推出 \(\mathbf{a}_1g_1 \ll g_2\) 与 \(\mathbf{f} - \mathbf{a}_1g_1 \ll g_1 \ll g_2\) ，故 \(\mathbf{f} \ll g_2\) ；但也将有 \(\mathbf{a}_2g_2 \ll g_2\) ，从而 \(\mathbf{a}_2 \ll 1\) ，与假设 \(\mathbf{a}_2 \neq 0\) 及 \((\mathrm{CO_{II}})\) 矛盾。所以必有 \(g_1 = g_2\) ；由关系 \(\mathbf{f} - \mathbf{a}_1g_1 \ll g_1, \mathbf{f} - \mathbf{a}_2g_1 \ll g_1\) 推出 \((\mathbf{a}_2 - \mathbf{a}_1)g_1 \ll g_1\) ，于是 \(\mathbf{a}_2 - \mathbf{a}_1 \ll 1\) ，从而 \(\mathbf{a}_2 = \mathbf{a}_1\) （由 \((\mathrm{CO_{II}})\) 与 \((\mathrm{CO_{III}})\) ）。

例. 当实数 \(x\) 趋于 \(+\infty\) 时， \(\mathbf{R}\) 上具有同一周期 \(\tau\) 的周期有界函数满足条件 \((\mathrm{CO}_{\mathrm{I}})\) 、 \((\mathrm{CO}_{\mathrm{II}})\) 与 \((\mathrm{CO}_{\mathrm{III}})\) ：若 \(\lim_{x\to +\infty}a(x) = 0\) ，则对每个 \(\varepsilon >0\) ，存在一个 \(x_0\) ，使得 \(|a(x)|\leqslant \varepsilon\) 对每个 \(x\geqslant x_0\) 成立；由此推出 \(|a(x)|\leqslant \varepsilon\) 对 \(0\leqslant x\leqslant \tau\) 也成立，因为存在一个整数 \(n\) 使得 \(x + n\tau \geqslant x_0\) ，且使得 \(a(x) = a(x + n\tau)\) ；由于 \(\varepsilon\) 是任意的，有 \(a(x) = 0\) 于 \([0,\tau]\) 上，从而处处为零。

用 V, p.~222 的记号，我们将说 \(\sum_{\lambda \leqslant \alpha} \mathbf{a}_{\lambda} g_{\lambda}\) （其中 \(\mathbf{a}_{\lambda}\) 属于 \(\mathcal{C}\) 且除有限多个外全为零）是 \(\mathbf{f}\) 的一个系数在 \(\mathcal{C}\) 中的、精确到 \(g_{\alpha}\) 的渐近展开，如果 \(\mathbf{f} - \sum_{\lambda \leqslant \alpha} \mathbf{a}_{\lambda} g_{\lambda} \ll g_{\alpha}\) ；于是对每个指标 \(\mu\) （满足 \(\mathbf{a}_{\mu} \neq 0\) ）， \(\mathbf{a}_{\mu} g_{\mu}\) 是 \(\mathbf{f} - \sum_{\lambda < \mu} \mathbf{a}_{\lambda} g_{\lambda}\) 的主部（相对于 \(\mathcal{E}\) 与 \(\mathcal{C}\) ），这证明了 \(\mathbf{f}\) 的渐近展开（精确到 \(g_{\alpha}\) ）存在时的唯一性。

在 \(n^{\circ}\) 3 (V, p.~223) 中给出的、构成 \(f_{1} + f_{2}\) 或 \([f_{1}.f_{2}]\) 的渐近展开（从 \(f_{1}\) 与 \(f_{2}\) 的给定渐近展开出发）的方法，仍可应用于变系数展开，只要 \([a_{\lambda}.b_{\mu}]\) 属于对应于赋范空间 V 的系数域 C，或容许一个系数在 C 中的渐近展开。

